

\section{Sobre el pegado de $\lambda$-celdas}

En esta sección sentaremos las bases de la Teoría de Morse. Hasta ahora hemos aportado el aparataje técnico sobre el que construir la teoría, aquí veremos cómo las funciones de Morse nos pueden servir para relacionar el tipo de homotopía de una variedad con el de un CW-complejo con tantas $\lambda$-celdas como puntos críticos de índice  $\lambda$ tenga $f$.

\medskip

Empezamos introduciendo algunos resultados sobre campos vectoriales (respetando el vocabulario de \cite{milnor1963}) que nos servirán para probar los resultados principales de la sección.

\iffalse
Antes de probar los resultados principales de esta sección conviene mencionar qué entendemos por \textit{variedad riemanniana} y dar una definición de campo vectorial.

\begin{definition}
Sea $M$ una variedad diferenciable.  
Llamamos \textit{fibrado tangente} de $M$ a la unión disjunta
de los espacios tangentes en todos los puntos de $M$, es decir,
\[
TM \;=\; \bigsqcup_{p \in M} T_pM.
\]
\end{definition}


\begin{definition}
Un \textit{campo vectorial suave} sobre una variedad diferenciable $M$ es una aplicación suave
\[
X \colon M \to TM,
\]
que denotamos $p \mapsto X_p$, y que verifica que, para todo $p \in M$, se tiene $X_p \in T_pM$.
\end{definition}


\begin{definition}
Una \textit{variedad riemanniana} es un par $(M,\langle\cdot,\cdot\rangle)$ donde $M$ es una variedad diferenciable y para cada punto $p \in M$, $\langle\cdot,\cdot\rangle_p$ es un producto escalar (luego, bilineal, simétrico y definido positivo)
  \[
    \langle\cdot,\cdot\rangle_p \colon T_pM \times T_pM \to \mathbb{R},
  \]
que verifica que para cualesquiera campos vectoriales suaves $X$ e $Y$, la función
    \[
      p \mapsto \langle X(p), Y(p) \rangle_p
    \]
    es suave sobre $M$.
\end{definition}
\fi

\begin{definition}
    Sea $M$ una variedad diferenciable. Llamamos \textit{grupo uniparamétrico de difeomorfismos de M} a una aplicación $\varphi:\mathbb{R}\times M \to M$ que verifica:
    \begin{enumerate}[label=\roman*)]
        \item $\forall t\in \mathbb{R},~~ \varphi_t:M\to M$ definida como $\varphi_t(q)=\varphi(t,q)$, es un difeomorfismo de $M$ sobre si misma.

        \item $\forall t, s\in\mathbb{R}$ se tiene que $\varphi_{t+s}=\varphi_t\circ \varphi_s$.
    \end{enumerate}
\end{definition}

\bigskip

Dado un grupo uniparamétrico de difeormorfismos $\varphi$ definimos el siguiente campo vectorial sobre $M$

\begin{equation*}
    X_q(f)=\lim_{h\to 0}\frac{f(\varphi_h(q))-f(q)}{h}
\end{equation*}
donde $f$ es una función suave de $M$ en $\mathbb{R}$, y decimos que $X$ genera $\varphi$.

Normalmente uno se refiere a la curva $\varphi_\bullet(q): \mathbb{R}\to M,~~ t\mapsto \varphi_t(q)$ como el \textit{flujo de X} cuando se conoce el campo vectorial $X$ que genera el grupo uniparamétrico de difeomorfismos. En cualquier caso, mantendremos un vocabulario lo más fiel posible al texto \cite{milnor1963}.

También hablaremos del \textit{vector velocidad} de una curva diferenciable $c:\mathbb{R}\to M$ que definimos de manera similar a X
\begin{equation*}
    \frac{dc}{dt}(f)= \lim_{h\to 0}\frac{f(c(t+h))-f(c(t))}{h}\in T_{c(t)}M 
\end{equation*} 

\bigskip
\begin{lemma}
\label{unidiff}
    Todo campo vectorial sobre una variedad diferenciable $M$ que se anula fuera de un compacto $K\subset M$ genera un único grupo de difeomorfismos de M.
\end{lemma}

\begin{proof}
    Sean $X$ un campo vectorial y $\varphi$ un grupo uniparamétrico de difeomorfismos de $M$ generado por $X$. Entonces se tiene que fijado un punto cualquiera $q\in M$ la curva $\varphi_\bullet(q): \mathbb{R}\to M,~~ t\mapsto \varphi_t(q)$ verifica la ecuación diferencial
    \begin{equation*}
        \frac{d\varphi_t}{dt}(q)=X_{\varphi_t(q)}\tag{$\ast$}
    \end{equation*}
    con condición inicial $\varphi_0(q)=q$.

    En efecto,
    \begin{align*}
        \frac{d\varphi_t(q)}{dt}(f)=\lim_{h\to 0} \frac{f(\varphi_{t+h}(q))- f(\varphi_t(q))}{h} = \lim_{h\to 0} \frac{f(\varphi_{h}(\varphi_t(q)))- f(\varphi_t(q))}{h} = \lim_{h\to 0} \frac{f(\varphi_{h}(p))- f(p)}{h} = X_p(f)
    \end{align*}

    con $p=\varphi_t(q)$.

    Ahora, por el Teorema de Picard-Lindelöf sabemos que existe localmente una solución única, que depende de manera diferenciable de la condición inicial. Es decir, $\forall p\in M$ existen un entorno $\mathcal{U}$ de $p$ y $\varepsilon > 0$ tales que $(\ast)$ tiene una única solución para todo $q\in \mathcal{U}$ y $|t|<\varepsilon$.

    Como $K\subset M$ es compacto, puede ser cubierto por una cantidad finita de dichos entornos. Sea $\varepsilon_0$ el menor de los valores definidos arriba y establezcamos $\varphi_t(q)=q,~~\forall q\notin K$ (luego $X_q=0$ fuera de $K$). Entonces $(\ast)$ tiene una única solución $\varphi_t(q)$ para $|t|<\varepsilon_0$ y $\forall q\in M$ suave respecto de ambas variables.

    Veamos que para todo $t$ adecuado cada $\varphi_t$ es un difeomorfismo de $M$ en si misma. Para ello basta comprobar que se cumple que $\varphi_{t+s}=\varphi_t\circ \varphi_s$ suponiendo que $|t|,|s|,|t+s|<\varepsilon_0$.

    Sean $q\in M$ fijo y  $p=\varphi_s(q)$, entonces $\varphi_t(p)=\varphi_t(\varphi_s(q))$. Definimos la curva $\gamma(u)=\varphi_{u+s}(q)$, así, $\gamma(0)=p$ y 
    
    \begin{equation*}
        \frac{d\gamma}{du}=\frac{d\varphi_{u+s}}{du}(q)=X_{\varphi_{u+s}(q)}=X_{\gamma (u)}
    \end{equation*} 
    
    Por lo que $\gamma$ es también una solución de ($\ast$) con condición inicial $\gamma(0)=p$.

    Sin embargo, tambien $\varphi_u(p)$ es solución con condición inicial $\varphi_0(p)=p$. Luego por unicidad de las soluciones
    \begin{equation*}
        \varphi_{u+s}(q)=\gamma(u)=\varphi_u(p)=\varphi_u(\varphi_s(q))=\varphi_u\circ \varphi_s(q)
    \end{equation*}
    Siempre que $|u|, |u+s|<\varepsilon_0$.

    Por último, necesitamos definir $\varphi_t$ para todo $t\in \mathbb{R}$.

    Sea $t\in \mathbb{R}$. Sean $k\in \mathbb{Z}$ y $r\in \mathbb{R}$ con $|r|\le \frac{\varepsilon_0}{2}$ los únicos valores tales que $t=k\frac{\varepsilon_0}{2}+r$. Entonces definimos

    \begin{equation*}
        \varphi_t=\underbrace{\varphi_{\frac{\varepsilon_0}{2}}\circ \dots \circ \varphi_{\frac{\varepsilon_0}{2}}}_{k~veces}\circ \varphi_r 
    \end{equation*}

    Es claro que $\varphi:\mathbb{R}\times M\to M$ es un difeomorfismo único para todo $t\in \mathbb{R}$.

\end{proof}

\begin{definition}
    Sean $(M,\langle\cdot,\cdot\rangle)$ una variedad riemanniana y $f:M\to\mathbb{R}$ suave. Se define el \textit{gradiente de f} como el campo vectorial $\nabla f$ que verifica 
    \begin{equation*}
        \langle X,\nabla f\rangle = X(f)
    \end{equation*}
    para cualquier otro campo vectorial $X$.
\end{definition}
\bigskip
% aqui deberia escribir algo para introducir lo siguiente.

\begin{definition}
    Sean $M$ una variedad diferenciable y $f:M\to\mathbb{R}$ suave definimos los \textit{conjuntos subnivel de f} como
    \begin{equation*}
        M^a := f^{-1}(-\infty,a] =\{p\in M: f(p)\le a\}
    \end{equation*}
\end{definition}

\begin{theorem}
\label{retractodef}
    Sea $f:M\to\mathbb{R}$ suave sobre una variedad riemanniana $(M,\langle\cdot,\cdot\rangle)$. Sean $a<b$ reales tales que $f^{-1}[a,b]$ es compacto y no contiene ningún punto crítico de $f$. Entonces $M^a$ es difeomorfo a $M^b$. Además, $M^a$ es retracto de deformación de $M^b$, luego tienen en el mismo tipo de homotopía.
\end{theorem}

\begin{proof}
    Sea $c:\mathbb{R}\to M$ una curva con vector velocidad $\frac{dc}{dt}$, entonces
    \begin{equation*}
        \langle \frac{dc}{dt}, \nabla f\rangle = \frac{dc}{dt}(f)=\frac{d(f\circ c)}{dt}
    \end{equation*}

    Definimos la aplicación suave $\rho : M\to \mathbb{R}$ como

    \begin{equation*}
        \rho(p):=
        \begin{cases}
            \frac{1}{\langle \nabla_pf, \nabla_pf\rangle}~~,~p\in f^{-1}[a,b]\\
            0 ~~~~~~~~~~~~~~,p\notin K
        \end{cases}
    \end{equation*}

    donde $K$ es un entorno compacto de $f^{-1}[a,b]$

    Ahora, sea el campo vectorial dado por $X_q=\rho(q)\nabla_pf$. Se tiene que $X_p$ se anula fuera de $K$ luego por el Lema \ref{unidiff} genera un único grupo uniparamétrico de difeomorfismos $\varphi_t$ de $M$.

    Fijemos $q\in M$ y consideremos $f\circ\varphi_{\bullet}(q),~t\mapsto f(\varphi_t(q))$. Si $\varphi_t(q)$ está en $f^{-1}[a,b]$ se tiene
    \begin{align*}
        \frac{df(\varphi_t(q))}{dt} = \left\langle \frac{d\varphi_t(q)}{dt},\nabla f\right\rangle = \langle X, \nabla f\rangle = \langle \rho~ \nabla f, \nabla f\rangle =\\ =\left\langle \frac{1}{\langle \nabla f,\nabla f\rangle}\nabla f,\nabla f\right\rangle=
        \frac{\langle \nabla f,\nabla f\rangle}{\langle \nabla f,\nabla f\rangle}= 1,
    \end{align*}

    por lo que $f(\varphi_t(q))$ es lineal con derivada $1$ siempre que $f(\varphi_t(q))$ esté en $[a,b]$, es decir, $f(\varphi_t(q))=f(q)+t$.

    Veamos ahora que $\varphi_{b-a}$ es un difeomorfismo de $M^a$ en $M^b$.

    Sea $q\in M$,
    \begin{equation*}
        \varphi_{b-a}(q)\in M^b \iff(\varphi_{b-a}(q))=f(q)+b-a\le b \iff f(q)-a\le 0 \iff q\in M^a,
    \end{equation*}
     y en concreto manda la frontera de $M^a$ a la frontera de $M^b$
     \begin{equation*}
         q\in \partial M^a \iff f(q)=a \implies f(\varphi_{b-a}(q))=b \iff b\in \partial M^b.
     \end{equation*}

     Para concluir la demostración falta comprobar que efectivamente $M^a$ es retractor de deformación de $M^b$. Para ello vamos a construir una homotopía relativa a $M^a$ de la identidad de $M^b$ a una retracción de $M^b$ a $M^a$.
     
    Sea $H(t,q):[0,1]\times M^b\to M^b$ definida como
    \begin{equation*}
        H(t,q)=
        \begin{cases}
            q~~~~~~~~~~~~~~~~~,~f(q)\le a\iff q\in M^a,\\
            \varphi_{t(a-f(q))}(q)~~,a\le f(q)\le b.
        \end{cases}
    \end{equation*}

    Entonces, $H(0,\cdot) = id_{M^b}$ ya que $\varphi_0(q)=q,~\forall q\in M^b$ y $H(1,\cdot)$ es una retracción de $M^b$ en $M^a$, comprobémoslo. Sea $q\in f^{-1}[a,b]$, se tiene
    \begin{equation*}
        H(1,q)=f(\varphi_{(a-f(q))}(q))= f(q)+(a-f(q))= a \iff \varphi_{(a-f(q))}(q)\in f^{-1}(a),
    \end{equation*}

    luego $H(1,H(1,\cdot))= id_{M^a}$.    

\end{proof}
\begin{remark}
    Nótese la importancia de que $f^{-1}[a,b]$ sea compacto. Si por ejemplo nuestra variedad es la esfera $\mathbb{S}^2$ embebida en $\mathbb{R}^3$, es claro que la función que asigna a cada punto de la esfera su coordenada $z$ es de Morse y verifica que $f^{-1}[a,b]$ es compacto puesto que es cerrado y la esfera es compacta. Ahora, si eliminamos de $\mathbb{S}^2$ un punto $p$ tal que $f(p)$ queda entre $f(a)$ y $f(b)$ se tiene que $f^{-1}[a,b]$ deja de ser compacto (pues deja de ser cerrado y estamos en un espacio Hausdorff) y claramente $M^a$ no es retracto de deformación de $M^b$. 
\end{remark}
\bigskip

El siguiente teorema es crucial para conocer la estructura de CW-complejo de una variedad. Este establece que el cambio provocado en el tipo de homotopía de los conjuntos subnivel de una función de Morse al cruzar un punto crítico equivale al de pegar una célula de dimensión el índice de dicho punto crítico.

\begin{theorem}
\label{pegadocelula}
     Sea $f:M\to\mathbb{R}$ suave sobre una variedad riemanniana $(M,\langle\cdot,\cdot\rangle)$ y $p\in M$ un punto crítico no degenerado de $f$ con índice $\lambda$. Supongamos que $f(p)=c$ y que  $\exists\varepsilon>0$ tal que $f^{-1}[c-\varepsilon,c+\varepsilon]$ es compacto y no contiene ningún punto crítico de $f$ a parte de $p$. Entonces, para todo $0<\varepsilon_0<\varepsilon$ suficientemente pequeño, $M^{c+\varepsilon}_f$ tiene el tipo de homotopía de $M^{c-\varepsilon}_f$ con una $\lambda$-celda pegada.
\end{theorem}

\begin{proof}
    Vamos a definir una función $F$ que sea menor que $f$ en un entorno $\mathcal{U}$ de $p$ con el fin de dejar a $p$ fuera de $F^{-1}[c-\varepsilon, c+\varepsilon]$ y poder aplicar a $F$ el Teorema \ref{retractodef}. Así, habremos demostrado que $F^{-1}(-\infty,c-\varepsilon]=M^{c-\varepsilon}_F$ es retracto de deformación de $M^{c+\varepsilon}_f$.
    
    Para finalizar la demostración probaremos que $M^{c-\varepsilon}_f$ con una $\lambda$-celda pegada es retracto de deformacion de $M^{c-\varepsilon}_F$ y por la transitividad de esta relación habremos terminado.

    \bigskip
%Preliminares------------------

    Sean una carta $(\mathcal{U}, \phi)$ de $M$ que contiene a $p$ y $(u^1,\dots , u^n)$ las coordenadas inducidas tal que $u^1(p)=\dots=u^{n}(p)=0$. Con el fin de simplificar las expresiones que siguen definimos las funciones $\xi:\mathcal{U}\to [0,\infty)$ y $\eta:\mathcal{U}\to [0,\infty)$ como
    \begin{gather*}
        \xi(q)=(u^1(q))^2+\dots +(u^{\lambda}(q))^2\\
        \eta(q)=(u^{\lambda +1}(q))^2+\dots +(u^{n}(q))^2
    \end{gather*}
    Aplicando el Lema de Morse (\ref{lemamorse}) a $f$ obtenemos
    \begin{equation*}
        f(q) = c - \xi(q) + \eta(q)
    \end{equation*}
    $\forall q\in \mathcal{U}$, donde $c=f(p).$

    Sea $0<\varepsilon_0\le\varepsilon$ suficientemente pequeño para que:
    \begin{enumerate}[label=\roman*)]
        \item $f^{-1}[c-\varepsilon_0,c+\varepsilon_0]$ sea compacto y no contenga ningún punto crítico de $f$ distinto de $p$.
        \item $\phi^{-1}(\mathcal{U})\subset \mathbb{R}^n$ contenga a la bola cerrada
        \begin{equation*}
            \phi\left(\overline{B(0, \sqrt{2\varepsilon_0})}\right)= \{q\in \mathcal{U}: \xi(q)+\eta(q)\le 2\varepsilon_0\}
        \end{equation*}
    \end{enumerate}

%Construcción de F----------------
    Construyamos $F:M\to \mathbb{R}$. Sea $\mu:\mathbb{R}\to\mathbb{R}$ una aplicación suave que verifica lo siguiente
    
    \begin{equation*}
        \begin{cases}
        \mu(0)>\varepsilon_0\\
        \mu(r)=0, r\ge 2\varepsilon_0\\
        -1<\mu'(r)\le 0, \forall r\in \mathbb{R}
        \end{cases}
    \end{equation*}
    entonces definimos $F$ como
    \vspace{1em}
    \begin{equation*}
        F(q)=
        \begin{cases}
        f(q)~~~~~~~~~~~~~~~~~~~~~~~~,~~q\notin \mathcal{U}\\
        f(q)- \mu(\xi(q)+2\eta(q)),~~q\in \mathcal{U}.
        \end{cases}
    \end{equation*}

    Es claro que $F$ esta bien definida ya que $\mu$ se anula en $\partial\mathcal{U}$. Además, es suave ya que en un entorno de $\partial\mathcal{U}$ se tiene que $\mu'$ es identicamente $0$, luego tanto aproximándonos desde dentro de $\mathcal{U}$ como desde fuera $d_qF=d_qf$.

    \bigskip
    
%Primera retraccion-------------------------------

    Ahora que tenemos $F$ debemos comprobar que nos encontramos en las hipótesis del Teorema \ref{retractodef}. Veamos primero que $M_F^{c+\varepsilon_0}=M_f^{c+\varepsilon_0}$, después que $F$ y $f$ poseen los mismos puntos críticos y que $F^{-1}[c-\varepsilon_0,c+\varepsilon_0]$ no contiene ningún punto crítico.

    En efecto, $\mu(\xi(q)+2\eta(q))=0, ~\forall q\in M$ tal que $\xi(q)+2\eta(q)\ge 2\varepsilon_0$ luego fuera del elipsoide $E=\{q\in M:\xi(q)+2\eta(q)\le 2\varepsilon_0\}$ se tiene que $F=f$.

    \medskip
    
    Sabemos que $E\in M_f^{c+\varepsilon_0}$ y por
    \begin{equation*}
        F(q)<f(q)=c-\xi(q)+\eta(q)\le c+\frac{1}{2}\xi(q)+\eta\le c + \varepsilon_0, ~\forall q\in E,
    \end{equation*}

    se tiene que $E\in M_F^{c+\varepsilon_0}$. Visto que es la única región en la que difieren concluimos que $M_F^{c+\varepsilon_0}=M_f^{c+\varepsilon_0}$.

    \medskip

    Ahora veamos que coinciden sus puntos críticos. Observamos que en $\mathcal{U}$ se verifican
    \begin{align*}
        \frac{\partial F}{\partial \xi}&=-1-\mu'(\xi + 2\eta)<0 ~\text{ y}\\
        \frac{\partial F}{\partial \eta}&=1-2\mu'(\xi+2\eta)\ge 1.
    \end{align*}
    Los puntos críticos de $F$ serán los que anulen su derivadas parciales
    \begin{align*}
        \frac{\partial F}{\partial x_i}=\frac{\partial F}{\partial \xi}\frac{\partial \xi}{\partial x_i} + \frac{\partial F}{\partial \eta}\frac{\partial \eta}{\partial x_i},
    \end{align*}

    pero el único punto que verifica simultáneamente $\frac{\partial \xi}{\partial x_i}=0$ y $\frac{\partial \eta}{\partial x_i}=0$ es $p$ ya que si no $f$ tendría más puntos críticos en $\mathcal{U}$. Luego como fuera de $\mathcal{U}$ coinciden se tiene que  $F$ y $f$ tienen los mismos puntos críticos.

    \medskip

    Veamos que $F^{-1}[c-\varepsilon_0,c+\varepsilon_0]=M^{c+\varepsilon_0}_F\setminus M^{c-\varepsilon_0}_F$ es compacto y no contiene ningún punto crítico. Es claro que $M^{c-\varepsilon_0}_f\subset M^{c-\varepsilon_0}_F$ ya que $F$ y $f$ coinciden salvo en el entorno de $p$ en el que $F\le f$ así
    \begin{equation}
        M^{c+\varepsilon_0}_F=M^{c+\varepsilon_0}_f\implies M^{c+\varepsilon_0}_F\setminus M^{c-\varepsilon_0}_F \subset M^{c+\varepsilon_0}_f\setminus M^{c-\varepsilon_0}_f=f^{-1}[c-\varepsilon_0,c+\varepsilon_0]
    \end{equation}
    por lo que $F^{-1}[c-\varepsilon_0,c+\varepsilon_0]$ es compacto. Por otro lado,
    \begin{equation*}
        F(p)=c-\mu(0)<c-\varepsilon_0
    \end{equation*}

    luego $p\notin F^{-1}[c-\varepsilon_0,c+\varepsilon_0]$.

    \medskip

    De esta manera, nos encontramos en las hipótesis del Teorema \ref{retractodef} para $F$, lo que demuestra que $M^{c-\varepsilon_0}_F$ es retracto de deformación de $M^{c+\varepsilon_0}_F=M^{c+\varepsilon_0}_f$.

    \bigskip

%Segunda retracción------------------

    Sea $e^{\lambda}=\{q\in \mathcal{U}: \xi(q)\le \varepsilon_0,~\eta(q)=0 \}\subset \mathcal{U}$ una $\lambda$-celda. En primer lugar, veamos que
    \begin{equation*}
        \partial e^{\lambda}=\{q\in \mathcal{U}: \xi(q)= \varepsilon_0,~~\eta(q)=0 \}=M^{c-\varepsilon_0}_f\cap e^{\lambda}.
    \end{equation*}

    \begin{quote}
    $\left( \subseteq \right) ~~\forall q\in \partial e^{\lambda},~f(q)=c-\varepsilon_0 \implies q\in M^{c-\varepsilon_0}_f$
    
    $\left( \supseteq \right) ~~\forall q\in M^{c-\varepsilon_0}_f\cap e^{\lambda}$, se tiene $\xi(q)\le\varepsilon_0,~\eta(q)=0 \text{ y } f(q)\le -\varepsilon_0$, entonces como
    \begin{equation*}
        f(q)=c-\xi(q)\le c-\varepsilon_0\iff \xi(q)\ge \varepsilon_0
    \end{equation*}
    \quad \quad necesariamente $\xi(q)=\varepsilon_0$, luego $q\in \partial e^{\lambda}$. 
    \end{quote}
    \bigskip

    De cara al resto de la demostración resultará conveniente introducir la noción de \textit{asa}. Entenderemos $M_F^{c-\varepsilon_0}$ como $M_f^{c-\varepsilon_0}\cup H$ donde $H=\overline{M_F^{c-\varepsilon_0}\setminus M_f^{c-\varepsilon_0}}$ es el asa del que hablamos, denotado así en virtud de su nombre en inglés \textit{handle}.

    \medskip
    
    Observamos que $e^{\lambda}$ está contenido en $H$. En efecto, $F$ alcanza su máximo en $\mathcal{U}$ cuando $\mu(\xi+2\eta)$ alcanza su mínimo y este lo alcanza cuando $\xi$ y $\eta$ se anulan, es decir, únicamente en $p$, así
    \begin{equation*}
        F(q)\le F(p)<c-\varepsilon_0 ~\text{ y }~f(q)=c-\xi(q)\ge c-\varepsilon_0,~\forall q\in e^{\lambda}
    \end{equation*}
    luego $q\in e^{\lambda}$ lo que implica que $q$ está en $H$.

    \medskip

    
    Ahora queremos demostrar que $M^{c-\varepsilon_0}_f\cup e^{\lambda}$ es retracto de deformación de $M^{c-\varepsilon_0}_F = M_f^{c-\varepsilon_0}\cup H$. Para ello vamos a construir una homotopía relativa a $M^{c-\varepsilon_0}_f\cup e^{\lambda}$ de la identidad a una retracción en $M^{c-\varepsilon_0}_f\cup e^{\lambda}$, como en la demostración del Teorema \ref{retractodef}.\\

    Sea $D:[0,1]\times M_f^{c-\varepsilon_0}\cup H\to M_f^{c-\varepsilon_0}\cup H$, definimos $D$ como la identidad fuera de $\mathcal{U}$. En $\mathcal{U}$ debemos distinguir tres regiones.\\

    \textbf{Región 1:}\quad $R_1=\{q\in\mathcal{U}:\xi(q)\le \varepsilon_0\}$.\\
    
    En esta región lo que queremos es llevar los puntos a $e^{\lambda}$.

    Definimos $D(t,q)=D(t,(u_1(q),\dots,u_n(q))=(u_1(q),\dots,u_{\lambda}(q),tu_{\lambda+1}(q),\dots,tu_n(q))$.\\
    
    Es claro que es relativa a $M^{c-\varepsilon_0}_f\cup e^{\lambda}$ ($R_1\cap M^{c-\varepsilon_0}_f=\varnothing$ y los puntos de $e^{\lambda}$ no los toca) y que  $D(1,\cdot)$ es la identidad. Además,
    \begin{equation*}
        \forall q\in R_1, ~D(0, q)\in e^{\lambda} ~\text{ y }~ D(0,D(0,q))=D(0,q).
    \end{equation*} 

    Asimismo, como $\frac{\partial F}{\partial \eta}>0$ (es decir, $F$ crece en la dirección de $\eta$) se tiene que $F(D(t,q))\le F(q)$ para todo $t\in [0,1]$ y todo $q\in R_1$. Luego $D$ lleva $M_f^{c-\varepsilon_0}\cup H$ en si mismo.
    \bigskip

    \textbf{Región 2:}\quad $R_2=\{q\in \mathcal{U}:~ \varepsilon_0\le \xi(q)\le \eta(q)+\varepsilon_0\}$.
    \bigskip

    Aquí lo que buscamos es llevar los puntos de $R_2$ a la frontera de $M_f^{c-\varepsilon_0}$.

    Definimos $D(t,q)=D(t,(u_1(q),\dots,u_n(q))=(u_1(q),\dots,u_{\lambda}(q),s_t(q) u_{\lambda+1}(q),\dots,s_t(q) u_n(q))$, donde 
    \begin{equation*}
        s_t(q)=t + (1-t)\left(\frac{\xi(q)-\varepsilon_0}{\eta}\right)^{1/2} \in ~[0,1]
    \end{equation*}

    Igual que sobre $R_1$ queda claro que es relativa a $M_f^{c-\varepsilon_0}\cup e^{\lambda}$, que $D(1,\cdot)$ es de nuevo la identidad y que $D(0,\cdot)$ es una retracción sobre la frontera de $M_f^{c-\varepsilon_0}$. Lo último se sigue de que en la frontera de $M_f^{c-\varepsilon_0}$ se tiene que $-\xi +\eta = \varepsilon_0$, luego $s_0(q)=1$. 
    
    Resta ver que $s_t$ es continua cuando $\xi \to \varepsilon_0$ y $\eta \to 0$. Reordenando la desigualdad que define $R_2$ tenemos $0\le\xi-\varepsilon_0\le\eta$, lo que significa que a medida que $\eta$ tiende a cero, $\xi$ debe tender a $\varepsilon_0$ más rápidamente. Así, 
    \begin{equation*}
        0\le \lim_{\xi \to \varepsilon_0, \eta\to 0} \frac{\xi - \varepsilon_0}{\eta} \le 1
    \end{equation*}
    y en particular, se tiene que $\lim_{\xi \to \varepsilon_0, \eta\to 0} \frac{\xi - \varepsilon_0}{\eta}=0$ necesariamente. Concluimos entonces que $s_t$ no genera problemas en la continuidad de $D$ cerca de los puntos con $\eta(q)=0$.
    \bigskip

    \textbf{Región 3:}\quad $R_3=\{q\in \mathcal{U}:~ \eta(q) + \varepsilon_0 \le \xi(q)\}$
    \bigskip

    $R_3$ coincide con $M_f^{c-\varepsilon_0}$ luego basta con tomar $D(t,q)= Id_{M_f^{c-\varepsilon_0}}$.
    \bigskip

    Las tres regiones cubren por completo $\mathcal{U}$. Esto concluye la demostración de que $M^{c-\varepsilon_0}_f\cup e^{\lambda}$ es retracto de deformación de $M^{c-\varepsilon_0}_F = M_f^{c-\varepsilon_0}\cup H$ y con ello la demostración del teorema.
    
\end{proof}
\begin{remark}
\label{obs_incluso_con_c_critico}
    Modificando un poco la demostración es posible demostrar que $M^c_f$ es también retracto de deformación de $M^{c+\varepsilon}_f$. De hecho, $M^c_f$ es retracto de deoformación de $M_F^c$, que es a su vez retracto de deformación de $M^{c+\varepsilon}_f$, esto, junto con el teorema, deja claro que $M^{c-\varepsilon}_f\cup e^{\lambda}$ es retracto de deformación de $M^c_f$.
\end{remark}

\bigskip

Uno puede demostrar de manera similar el caso general:

\begin{theorem}
\label{pegadocelula_general}
    Sean $f$ y $M$ en las hipótesis del teorema anterior. Sean $p_1,\dots, p_k$ puntos críticos no degenerados de $f$ con índices $\lambda_1,\dots, \lambda_k$ tales que $f(p_i)=c, ~\forall i\in \{1,\dots, k\}$. Entonces existe un $\varepsilon>0$ lo suficientemente pequeño tal que $M_f^{c+\varepsilon}$  tiene el tipo de homotopía de 
    \begin{equation*}
        M_f^{c-\varepsilon}\cup e^{\lambda_1}\cup\dots\cup e^{\lambda_k}.
    \end{equation*}
\end{theorem}
\bigskip


El último resultado corona la sección relacionando el tipo de homotopía de una variedad diferenciable con el de un CW-complejo cuyas celdas vienen definidas por los puntos críticos de una función de Morse. 

Su demostración necesitará del \textit{teorema de aproximación celular} cuya demostración se puede encontrar en \cite[p.~349]{hatcher2002}, del teorema $1$ de \cite{Whitehead1949CombinatorialHomotopy} y de dos lemas acerca de un espacio topológico con una celda pegada que probaremos. Procedemos en este orden.
\vspace{0.5em}

\begin{definition}
    Se dice que una función $f:X\to Y$ entre CW-complejos es una \textit{función celular} si $f(X^k)\subset Y^k$ para todo $k$. Es decir, una función que lleva celdas a celdas de su misma dimensión o más pequeña.
\end{definition}
\begin{theorem}[Teorema de aproximación celular]
\label{aprox_celular}
    Toda función $f:X\to Y$ entre CW-complejos es homótopa a una función celular.
\end{theorem}

\bigskip

\begin{definition}
    Se dice que un espacio topológico $X$ \textit{domina} a un espacio topológico $Y$ si existen funciones continuas $f:Y\to X$ y $g:X\to Y$ tales que $f\circ g$ es homótopa a la identidad de $X$.
\end{definition}

\begin{definition}
    Sea $X$ un espacio dominado por un CW-complejo $K$. Denotamos por $\Delta X$ la mínima dimensión de entre los CW-complejos que dominan a $X$ (puede ser infinita), es decir,
    \begin{equation*}
        \Delta X := \text{min}\{\text{dim}K\colon K \text{ es un CW-complejo que domina a }X\}.
    \end{equation*}
\end{definition}

El siguiente teorema requiere de más vocabulario aún para ser entendido. Presentamos aquí los objetos pero instamos al lector a recurrir a cualquiera de las referencias \cite{Whitehead1949CombinatorialHomotopy}, \cite{munoz2021} y \cite{hatcher2002} para mayor profundidad.

Dado un espacio topológico $X$ sus grupos de homotopía $\pi_n(X)$ son una generalización en dimensión superior de su grupo fundamental $\pi_1(X)$ formado por las clases de equivalencia módulo homotopía de lazos en $X$. Toda función continua $f:X\to Y$ induce homomorfismos $f_n:\pi_n(X)\to \pi_n(Y)$ entre dichos grupos de homotopía, los cuales serán isomorfismos si y solo si $f$ es una equivalencia homotópica.

\begin{theorem}
\label{dominados}
    Sean $X$ e $Y$ dos espacios topológicos conexos dominados por sendos CW-complejos. Entonces una función $f:X\to Y$ es una equivalencia homotópica si y solo si $f_n:\pi_n(X)\to \pi_n(Y)$ es un isomorfismo de grupos para todo $n$ tal que $1\le n<N+1$, donde $N=\text{max } \{\Delta X,\Delta Y\}$.
\end{theorem}

\bigskip

Vamos ahora con los dos lemas, ambos con demostración. Para la demostración del primero nos hemos basado en la demostración del lema $5$ encontrado en \cite[p.~67]{whitehead_lemacw}.

\begin{lemma}[Lema de Whitehead]
\label{lema1_democwvariedad}
    Sean $\varphi_0, \varphi_1: \partial e^{\lambda} \to X$ dos funciones de pegado homotópicas. Entonces, la función identidad de $X$ se puede extender a una equivalencia homotópica
    \begin{equation*}
        k: X\cup_{\varphi_0} e^{\lambda} \to X\cup_{\varphi_1} e^{\lambda}.
    \end{equation*}
\end{lemma}
\begin{proof}

    Primero definimos la función $k$ y su inversa $l$ y después probamos que, en efecto, dan una equivalencia homotópica. 
    
    Sea $\{\varphi_t:\partial e^{\lambda} \to X\}$ una homotopía entre $\varphi_0$ y $\varphi_1$, definimos $k$ como sigue.
    \begin{equation*}
        \begin{cases}
            k(x) = x &~~\text{ si }~~ x\in X\\
            k(tu) = 2tu &~~\text{ si }~~ 0\le t\le \frac{1}{2},~u\in \partial e^{\lambda}\\
            k(tu) = \varphi_{2-2t}(u) &~~\text{ si }~~  \frac{1}{2} \le t\le 1,~u\in \partial e^{\lambda},
        \end{cases}
    \end{equation*}
    donde $tu$ denota el producto usual de el escalar $t\in [0,1]$ con el vector unitario $u\in \partial e^{\lambda}$, es decir, cubre todo $D^{\lambda}$. Es continua ya que para $t=\frac{1}{2}$ se tiene
    \begin{equation*}
        k(u/2) = u = \varphi_{1}(u) ~~\text{ y }~~ u \sim \varphi_{1}(u) ~\text{ en }~ X\cup_{\varphi_1} e^{\lambda}.
    \end{equation*}
    Además, está bien definida ya que envía representantes de la misma clase de euivalencia al mismo punto. En $X\cup_{\varphi_0} e^{\lambda}$ están relacionado $u$ y $\varphi_0(u)$ para todo $u$ en $\partial e^{\lambda}$ y $k$ verifica que 
    \begin{equation*}
        k(u) = \varphi_{0}(u) ~\text{ y }~ k(\varphi_{0}(u)) = \varphi_{0}(u),
    \end{equation*} 
    donde la segunda igualdad se tiene ya que $\varphi_{0}(u)$ está en $X$.
      
    De manera similar, construimos $l: X\cup_{\varphi_1} e^{\lambda} \to X\cup_{\varphi_0} e^{\lambda}$ como
    \begin{equation*}
        \begin{cases}
            l(x) = x &~~\text{ si }~~ x\in X,\\
            l(tu) = 2tu &~~\text{ si }~~ 0\le t\le \frac{1}{2},~u\in \partial e^{\lambda},\\
            l(tu) = \varphi_{2t-1}(u) &~~\text{ si }~~  \frac{1}{2} \le t\le 1,~u\in \partial e^{\lambda}.
        \end{cases}
    \end{equation*}
    Se prueba de manera análoga a $k$ que $l$ es continua y está bien definida.
    
    Ahora, veamos como es $l\circ k$ y demos una homotopía de esta a la identidad de $X\cup_{\varphi_0} e^{\lambda}$. El caso $k\circ l$ es análogo y se deja a la discreción del lector.
    \begin{equation*}
        \begin{cases}
            l\circ k(x) = x &~~\text{ si }~~ x\in X,\\
            l\circ k(tu) = 4tu &~~\text{ si }~~ 0\le t\le \frac{1}{4},~u\in \partial e^{\lambda},\\
            l\circ k(tu) = \varphi_{4t-1}(u) &~~\text{ si }~~  \frac{1}{4} \le t\le \frac{1}{2},~u\in \partial e^{\lambda},\\
            l\circ k(tu) = \varphi_{2-2t}(u) &~~\text{ si }~~  \frac{1}{2} \le t\le 1,~u\in \partial e^{\lambda}.
        \end{cases}
    \end{equation*}
    El último caso se sigue de que $k(tu) = \varphi_{2-2t}(u)$ si $\frac{1}{2} \le t\le 1$ y esto pertecene a $X$, luego $l$ lo deja fijo. Veamos ahora la homotopía, definimos
    \begin{equation*}
    \begin{cases}
        h_s(x) = x &~~\text{ si }~~ x\in X, ~~s\in[0,1],\\[4pt]
        h_s(tu) = (4-3s)tu &~~\text{ si }~~ 0 \le t \le \frac{1}{4-3s},~u\in \partial e^{\lambda} \\[4pt]
        h_s(tu) = \varphi_{(4-3s)t-1}(u) &~~\text{ si }~~ \frac{1}{4-3s} \le t \le \frac{2-s}{4-3s},~u\in \partial e^{\lambda} \\[4pt]
        h_s(tu) = \varphi_{\frac{1}{2}(4-3s)(1-t)}(u),
        &~~\text{ si }~~  \frac{2-s}{4-3s} \le t \le 1,~u\in \partial e^{\lambda}.
    \end{cases}
    \end{equation*}
    Es claro que para $s=0$ recuperamos $l\circ k$ y para $s=1$ se tiene la identidad de $X\cup_{\varphi_0} e^{\lambda}$. La continuidad se demuestra fácilmente viendo el comportamiento en los extremos de cada tramo. Que $h_s$ está bien definida se comprueba como en el caso de $k$.

    Así concluye la demostración.

\end{proof}
\bigskip

\begin{lemma}
\label{lema2_democwvariedad}
    Sea $\varphi: \partial e^{\lambda} \to X$ una función de pegado. Toda equivalencia homotópica $f:X\to Y$ se puede extender a una equivalencia homotópica 
    \begin{equation*}
        F: X\cup_{\varphi} e^{\lambda} \to Y\cup_{f\circ\varphi} e^{\lambda}.
    \end{equation*}
\end{lemma}
\begin{proof}
    Definimos $F$ mediante
    \begin{equation*}
        \begin{cases}
            F|_X = f, \\
            F|_{e^{\lambda}} = \text{id}_{e^{\lambda}}.
        \end{cases}
    \end{equation*}

    Sea $g:Y\to X$ una inversa homotópica de $f$, definimos
    \begin{equation*}
        G: Y\cup_{f\circ\varphi} e^{\lambda} \to X\cup_{g\circ f\circ\varphi} e^{\lambda}
    \end{equation*}
    de la misma forma que $F$,
    \begin{equation*}
        \begin{cases}
            G|_Y = g, \\
            F|_{e^{\lambda}} = \text{id}_{e^{\lambda}}.
        \end{cases}
    \end{equation*}

    Ahora, como $g\circ f$ es homótopa a la identidad de $X$ tenemos que $g\circ f\circ\varphi$ es homótopa a $\varphi$. Luego se sigue del lema anterior que existe una equivalencia homotópica 
    \begin{equation*}
        k: X\cup_{g\circ f\circ\varphi} e^{\lambda} \to X\cup_{\varphi} e^{\lambda}.
    \end{equation*}
    que construimos igual que en la demostración de dicho lema usando una homotopía de $g\circ f$ a la identidad de $X$ $\{h_t: X \to X\}$.

    Veamos primero que la composición $(k\circ G\circ F):X\cup_{\varphi} e^{\lambda}\to X\cup_{\varphi} e^{\lambda}$ es homótopa a la identidad, de ahora en adelante nos referiremos a esta composición como $kGF$. 
    
    Es sencillo comprobar que $kGF$ viene dada por
    \begin{equation*}
        \begin{cases}
            kGF(x) = g(f(x)) &~~\text{ si }~~ x\in X,\\
            kGF(tu) = 2tu &~~\text{ si }~~ 0\le t\le \frac{1}{2},~u\in \partial e^{\lambda},\\
            kGF(tu) = h_{2-2t}(\varphi(u)) &~~\text{ si }~~ \frac{1}{2}\le t\le 1,~u\in \partial e^{\lambda}.
        \end{cases}
    \end{equation*} 
    Ahora, denotamos por $q_s:X\cup_{\varphi} e^{\lambda}\to X\cup_{\varphi} e^{\lambda}$ la homotopía que necesitamos de $kGF$ a la identidad. Esta viene dada por
    \begin{equation*}
        \begin{cases}
            q_s(x) = h_s(x) &~~\text{ si }~~ x\in X,\\
            q_s(tu) = \frac{2}{1+s}tu &~~\text{ si }~~ 0\le t\le \frac{1+s}{2},~u\in \partial e^{\lambda},\\
            q_s(tu) = h_{2-2t+s}(\varphi(u)) &~~\text{ si }~~ \frac{1+s}{2}\le t\le 1,~u\in \partial e^{\lambda}.
        \end{cases}
    \end{equation*}

    En $s=0$ recuperamos $kGF$ ya que $h_0\equiv g\circ f$, los otros dos tramos son claros. En $s=1$ se tienen $h_1$, que es la identidad en $X$,  $q_1(tu)=tu$ para todo $t\in[0,1]$ y $q_1(u)=h_1(\varphi(u))$ en el último tramo (ya que $t$ solo puede tomar el valor $1$), que es igual a $\varphi(u)$ que está relacionado con $u$. Se comprueba fácilmente que no hay problemas de continuidad observando los valores que toma $q$ cuando $t=\frac{1+s}{2}$ están relacionados.

    Además, $q$ está bien definida ya que $q_s(u) = h_s(\varphi(u)) = q_s(\varphi(u))$.

    Así hemos probado que $F$ tiene una inversa homotópica por la izquierda. Se prubea de manera análoga que $G$ tiene un inversa homotópica por la derecha. 
    
    Ahora debemos probar que $kG$ es, en efecto, también una inversa homotópica por la derecha de $F$. Para ello probamos el siguiente resultado primero:
    \begin{quote}
        Si una función $F:X\to Y$ tiene una inversa homotópica por la izquierda $L$ y una inversa homotópica por la derecha $R$, entonces $F$ es una equivalencia homotópica y tanto $R$ como $L$ son inversas homotópicas por ambos lados.
    \end{quote}
    \medskip
    En efecto, las relaciones $L\circ F \simeq \text{id}_{X}$ y $F\circ R \simeq \text{id}_{Y}$ implican  
    \begin{equation*}
        L \simeq L(F\circ R) = (L\circ F)(R) \simeq R.
    \end{equation*}
    Luego, como $L$ esta relacionado con $R$,
    \begin{equation*}
        L\circ F \simeq R\circ F \simeq \text{id}_{X}.
    \end{equation*}
    Así, $R$ es una inversa homotópica por ambos lados. Se prueba para $L$ de la misma forma.

    \bigskip

    Terminamos la prueba del lema como sigue. Sabemos que tanto $F$ como $G$ tienen una inversa homotópica por la izquierda. 

    Ahora, como $k(GF)\simeq \text{id}_{X\cup_{\varphi} e^{\lambda}}$ y sabemos que $k$ tiene una inversa homotópica por la izquierda, se tiene que $(GF)k\simeq \text{id}_{X\cup_{g\circ f\circ\varphi} e^{\lambda}}$. Donde hemos aplicado el desarrollo dela demostración del resultado intermedio con $l=L$, $F=k$ y $R=GF$.

    De la misma forma, como $G(Fk) \simeq \text{id}_{X\cup_{g\circ f\circ\varphi} e^{\lambda}}$ y sabemos que $G$ tiene una invera homotópica por la izquierda, se tiene $(Fk)G \simeq \text{id}_{Y\cup_{f\circ\varphi} e^{\lambda}}$.

    Y, por último, como $F(kG) \simeq \text{id}_{Y\cup_{f\circ\varphi} e^{\lambda}}$ y $F$ tiene una inversa homotópica por la izquierda, se tiene que $F$ es una equivalencia homotópica, lo que conlcluye la demostración.

\end{proof}

Demostrados ambos lemas, procedemos a enunciar y demostrar el último resultado de la sección. 

\begin{theorem}
\label{cwvariedad}
    Si $f$ es una función de Morse sobre una variedad diferenciable $M$ y $M^a$ es compacto para cualquier $a\in\mathbb{R}$. Entonces $M$ tiene el tipo de homotopía de un CW-complejo con una $\lambda$-celda por cada punto crítico de índice $\lambda$.
\end{theorem}
\begin{proof}
    Sea $\{c_i\}_{i\in \mathcal{C}}$ el conjunto de los valores críticos de $f$ ordenados de tal manera que se verifica $c_1<c_2<_3<\dots$ Este conjunto no puede tener ningún punto de acumulación, para verlo, supongamos que tuviera uno.

    En primer lugar, que el conjunto $\{c_i\}_{i\in \mathcal{C}}$ tenga un punto de acumulación $c$ implica que existe una sucesión $\{c_{i_k}\}_{k\in \mathbb{N}}$ convergente a $c$. Ahora, por esto, existen inifnitos $c_i$ menores que $c+1$ y como, por definición, para cada $c_i$ existe un punto crítico $p_i\in M$ tal que $f(p_i) = c_i$ tenemos que $M^{c+1}$ contiene infinitos puntos críticos de $f$. 
    
    Esta colección infinita de puntos críticos contiene necesariamente una sucesión $\{p_{i_k}\}_{k\in \mathbb{N}}$ convergente a un punto $p$ de $M^{c+1}$. Por continuidad de $f$ se tiene que 
    \begin{equation*}
        f(p) = \lim_{k\to \infty} f(p_{i_k}) =  \lim_{k\to \infty} c_{i_k} = c
    \end{equation*}
    y por continuidad de la diferencial de $f$, $p$ es un punto crítico. Esto es una contradicción ya que $f$ es una función de Morse, luego $p$ es no degenerado y, por tanto, aislado.

    Sea $a\in\mathbb{R}\setminus\{c_i\}_{i\in \mathcal{C}}$, procedemos por inducción completa en la cantidad de valores críticos menores que $a$. 
    
    Para el caso base supongamos $a<c_1$. Es claro que $M^a$ es vacío ya que si no lo fuera, por ser compacto, $f$ debe alcanzar un mínimo en $M^a$ y tendría un valor crítico menor que $c_1$.

    Para la hipótesis de inducción supongamos que $M^a$ con $c_{i-1}<a<c_{i}$ tiene el tipo de homotopía de un CW-complejo $K$ y denotemos por $h:M^a\to K$ una equivalencia homotópica. Veamos entonces que dado $a'>c_{i}$ se tiene que $M^{a'}$ tiene el tipo de homotopía de un CW-complejo. 
    
    Por el teorema \ref{pegadocelula_general}, para $\varepsilon$ suficientemente pequeño $M^{c_i+\varepsilon}$ tiene el tipo de homotopía de 
    \begin{equation*}
        M^{c_i - \varepsilon}\cup_{\varphi_1} e^{\lambda_1} \cup\dots\cup_{\varphi_m} e^{\lambda_m}
    \end{equation*}
    para ciertas funciones de pegado $\varphi_1:\partial D^{\lambda_1}\to M^{c_i - \varepsilon},\dots,\varphi_m: \partial D^{\lambda_m}\to M^{c_i - \varepsilon}$, tomemos $a'=c_i+\varepsilon$. Además, por el teorema \ref{retractodef} existe una equivalencia homotópica $g: M^{c_i-\varepsilon}\to M^a$.

    Ahora, por el teorema de aproximación celular (\ref{aprox_celular}) se tiene que para todo $j\in\{1,\dots,m\}$ $h\circ g\circ\varphi_j$ es homótopa a una función
    \begin{equation*}
        \psi_j: \partial D^{\lambda_j} \to K^{\lambda_j - 1},
    \end{equation*}
    donde $K^{\lambda_j - 1}$ denota el $(\lambda_j -1)$-esqueleto de $K$. Así, por los lemas \ref{lema1_democwvariedad} y \ref{lema2_democwvariedad} se tiene que 
    \begin{equation*}
        K\cup_{\psi_1} e^{\lambda_1} \cup\dots\cup_{\psi_m} e^{\lambda_m}
    \end{equation*}
    es un CW-complejo y tiene el mismo tipo de homotopía que $M^{c_i+\varepsilon}$. Hemos completado la inducción, $M^a$ tiene el tipo de homotopía de un CW-complejo para todo $a$.

    Si $M$ es compacta hemos terminado ya que tendríamos una cantidad finita de valores críticos y se tiene que $M$ es igual a $M^a$ para todo $a\ge \text{max }f$, es decir, la inducción cubre todo $M$.

    Si $M$ no es compacta pero la cantidad de valores críticos es finita, se tiene que todos los puntos críticos de $M$ se encuentran en un mismo $M^a$ y de manera similar a la demostración del teorema \ref{retractodef} se puede demostrar que $M^a$ es retracto de deformación de $M$ y, de nuevo, hemos terminado.

    Si tenemos una cantidad infinitad de valores críticos y, por tanto, de puntos críticos, la construcción hecha en la inducción nos da la siguiente secuencia de equivalencias homotópicas
    \[
        \begin{array}{cccccc}
            M^{a_1} & \subset & M^{a_2} & \subset & M^{a_3} & \subset \cdots \\
            \;\;\downarrow{\simeq} & & \;\;\downarrow{\simeq} & & \;\;\downarrow{\simeq} & \\
            K_1 & \subset & K_2 & \subset & K_3 & \subset \cdots
        \end{array}
    \]
    tales que cada una extiende a la anterior. Sea $K$ la unión de todos los $K_i$. Sea $\Phi:M\to K$ la función límite, se tiene que $\Phi$ induce isomorfismos entre los grupos de homotopía en todas las dimensiones, entonces, por el teorema \ref{dominados} es una equivalencia homotópica, lo que concluye la demostración.

\end{proof}
\begin{remark}
    Durante la demostración en la que nos basamos \cite[p.~23]{milnor1963}, Milnor omite algunos detalles técnicos. Aquí hemos tratado salvar todas las lagunas técnicas posibles, sin embargo, hay al menos dos que hemos obviado:
    \begin{enumerate}[label=(\roman*)]
        \item Al final de la demostración, para ver que se puede aplicar el teorema \ref{dominados} hay que comprobar que tanto $K$ como $M$ están dominados por CW-complejos. Es claro que $K$ está dominado por si mismo, para demostrar que $M$ lo está Milnor mencionar que ha de considerarse como un retacto de un entorno tubular del espacio euclideo.
        \item Asumimos que la función límite $\Phi$ existe sin dar demostración.
    \end{enumerate}    
\end{remark}

\begin{remark}
    Hemos probado que cada $M^a$ tiene el tipo de homotopía de un CW-complejo con una cantidad finita de celdas, una por cada punto crítico. Esto es cierto incluso cuando $a$ es un valor crítico (ver la observación \ref{obs_incluso_con_c_critico}).
\end{remark}

%%%%%%%%%%%%%%%%%%%%%%%%%%%%%%%
%%                           %%
%%       APLICACIONES        %%
%%                           %%
%%%%%%%%%%%%%%%%%%%%%%%%%%%%%%%
\subsection{Algunas aplicaciones}
\label{subsec:aplicaciones}
Presentamos brevemente dos ejemplos de aplicaciones de los teoremas acerca del pegado de $\lambda$-celdas.

\begin{theorem}[Teorema de Reeb]
    Si $M$ es una variedad diferenciable compacta y $f$ es una función de Morse sobre $M$ con solo dos puntos críticos, entonces $M$ es homeomorfa a una esfera.
\end{theorem}
\begin{proof}
    Dichos puntos críticos han de ser necesariamente el máximo y el mínimo. Llamemos $p$ y $q$ a los puntos críticos y digamos que $f(p)=0$ es el mínimo y $f(q)=1$ es el máximo.

    Sabemos por el lema de Morse (\ref{lemamorse}) que existen entornos $\mathcal{U}_p$ de $p$ y $\mathcal{U}_q$ de $q$ y sistemas de coordenadas locales $(x_1,\dots,x_n)$ y $(y_1,\dots,y_n)$ tales que $p$ y $q$ tienen, respectivamente, coordenadas $0$ y $f$ es de la forma
    \begin{equation*}
        f(x) = x_1^2+x_2^2+\dots+x_n^2~~\text{ y }~~ f(x)=1-y_1^2- y_2^2-\dots-y_n^2
    \end{equation*}
    Ahora, sea $\varepsilon > 0$ suficientemente pequeño, se tiene que
    \begin{equation*}
        M^{\varepsilon} = \{p\in M \colon f(p)\le \varepsilon\} ~~\text{ y }~~ f^{-1}[1-\varepsilon,1] = \{p\in M \colon 1-\varepsilon \le f(p)\le 1\}.
    \end{equation*}
    Ambos, en coordenadas locales se corresponden con $n$-discos cerrados de radio $\sqrt{\varepsilon}$ y centro $0$, por tanto, son homeomorfos a dichas bolas.

    Por otro lado, como $f^{-1}[\varepsilon, 1-\varepsilon]$ no contiene ningún punto crítico, por el teorema \ref{retractodef}$M^{\varepsilon}$ y $M^{1-\varepsilon}$ son difeomorfos. Así, $M$ es igual a la unión de los dos $n$-discos $M^{1-\varepsilon}$ y $f^{-1}[1-\varepsilon,1]$ pegados a lo largo de su frontera común.

\end{proof}
\vspace{-1em}
\begin{remark}
    \hfil
    \begin{enumerate}[label=(\roman*)]
        \item El resultado es cierto incluso cuando los puntos críticos son degenerados, sin embargo, la prueba es más complicada.
        \item Utilizando el teorema \ref{cwvariedad} es fácil ver que $M$ tiene el tipo de homotopía de un CW-complejo con una $0$-celda y una $n$-celda. Una esfera admite una estructura de CW-complejo con esas celdas, sin embargo, el teorema no nos dice como se pegan. Luego no podemos asegurar que $M$ tenga el tipo de homotopía de la esfera, lo que sería suficiente gracias a la \textit{conjetura de Poincaré generalizada}.
        \item No es necesario que $M$ sea difeomorfa a la esfera, esto se debe a la existencia de \textit{esferas exóticas} en algunas dimensiones, espacios homeomorfos a la esfera cuyas estructuras diferenciables no son difeomorfas a la de la esfera estándar.
    \end{enumerate}
\end{remark}

\vspace{2.5em}

Otra aplicación de los teoremas de la sección anterior es la siguiente.

Sea $\mathbb{CP}^n$ el espacio proyectivo complejo de dimensión $n$. Pensaremos en $\mathbb{CP}^n$ como el espacio cuyos puntos son las clases de equivalencia de $(n+1)$-uplas $(z_0,\dots,z_n)$ de números complejos denotadas por $[z_0:z_1:\dots:z_n]$ tales que $\sum_j |z_j|^2=1$. 

Definimos la función $f:\mathbb{CP}^n \to \mathbb{R}$ dada por 
\begin{equation*}
    f[z_0:z_1:\dots:z_n] = \sum_{j=0}^n c_j |z_j|^2,
\end{equation*}
donde $c_0,\dots,c_n\in \mathbb{R}$ son constantes distintas dos a dos. Para hallar sus puntos críticos consideremos el siguiente sistema de coordenadas locales de $\mathbb{CP}^n$. Sea $\mathcal{U}_0$ el subespacio de los puntos tales que $z_0$ es distinto de $0$, establecemos 
\begin{equation*}
    |z_0|\frac{z_j}{z_0} = x_j + iy_j
\end{equation*}
donde $x_1,y_1,\dots,x_n,y_n : \mathcal{U}_0\to \mathbb{R}$ son funciones coordenadas que llevan de manera difeomorfa $\mathcal{U}_0$ a la bola unidad de $\mathbb{R}^{2n}$. Es claro que 
\begin{equation*}
    |z_j|^2 = x_j^2 + y_j^2 ~~\text{ y }~~ |z_0|^2 = 1 - \sum_{j=1}^n (x_j^2 + y_j^2),
\end{equation*}
donde la segunda igualdad se sigue de que $\sum_j |z_j|^2=1$ implica que $|z_0|^2 = 1 - \sum_{j=1}^n |z_j|^2 $. Y por tanto $f$ en estas coordenadas locales es de la forma
\begin{equation*}
    f = c_0 + \sum_{j=1}^n (c_j-c_0)(x_j^2 + y_j^2)
\end{equation*}
ya que 
\begin{equation*}
    \sum_{j=0}^n c_j |z_j|^2 = c_0|z_0|^2 + \sum_{j=1}^n c_j |z_j|^2 = c_0 \Big(1 - \sum_{j=1}^n (x_j^2 + y_j^2) \Big) + \sum_{j=1}^n c_j(x_j^2 + y_j^2).
\end{equation*}
Así, obtenemos que para todo $j\in\{1,\dots,n\}$
\begin{equation*}
    \frac{\partial f}{\partial x_j} = 2x_j(c_j - c_0) ~~\text{ y }~~ \frac{\partial f}{\partial y_j} = 2y_j(c_j - c_0),
\end{equation*}
por lo que el único punto de $\mathcal{U}_0$ que anula todas las derivadas parciales es $p_0 = [1:0:\dots:0]$. 

Dado que la matriz Hessiana de $f$ en coordenadas es una matriz diagonal cuyos elementos son constantes no nulas, $p_0$ es no degenerado. Además, es sencillo ver que su índice es dos veces el número de constantes $c_j$ que son menores que $c_0$.

De manera análoga podemos considerar los sistemas de coordenadas equivalentes para los conjuntos $\mathcal{U}_1, \dots, \mathcal{U}_n$, obteniendo los puntos críticos no degenerados
\begin{equation*}
    p_1 = [0:1:0:\dots:0], \dots, p_1 = [0:0:\dots:0:1]
\end{equation*}
tales que el índice de $p_k$ es igual a dos veces el número de constantes $c_j$ menores que $c_k$. Como la unión de los $\mathcal{U}_j$ conforma el total de $\mathbb{CP}^n$ se sigue que estos son los únicos puntos críticos de $f$, por lo que todo índice posible entre $0$ y $2n$, es decir, los pares, se da una única vez.

Por el teorema \ref{cwvariedad} $\mathbb{CP}^n$ tiene el tipo de homotopía de un CW-complejo de la forma
\begin{equation*}
    e^0\cup e^2 \cup e^4 \cup \dots \cup e^{2n}.
\end{equation*}
